\FloatBarrier
\section{Measuring Interdot Tunneling}
\label{app:tunnel-coupling}
\label{sec:measurement-tunnel-spectroscopy}

The exchange result in Section~\ref{sec:exchange-virtual-tunneling}
depends on the coherent interdot tunnel coupling. This appendix shows
how to determine that parameter from the charge distribution of one
electron shared between two dots.

Consider one electron confined to two neighboring dots. Let
$\epsilon_1$ and $\epsilon_2$ be the energies of the states
localized in the left and right dots. Changing the gate voltages
varies their energy difference $\varepsilon=\epsilon_1-\epsilon_2$,
called detuning. We follow $n_L$, the probability of finding
the electron in the left dot, as a function of $\varepsilon$.
First consider the ground state at zero temperature. Without
tunneling, the electron occupies whichever dot has lower energy:
\[
 n_L(\varepsilon)=
 \begin{cases}
  1,&\varepsilon<0,\\
  0,&\varepsilon>0.
 \end{cases}
\]

Tunneling couples these two localized states. Let
$t_c=|t_{12}|$ be the magnitude of its matrix element, in energy
units. A large energy difference favors localization in the
lower-energy dot, while for $|\varepsilon|$ comparable to $t_c$
the ground state contains substantial contributions from both
dots. At a fixed nonzero $t_c$, the curve of $n_L$ against
$\varepsilon$ therefore falls continuously from near 1 to near 0,
instead of making a step at zero. A larger $t_c$ sustains this
mixing at a larger energy difference, so the fall spans a wider
interval of $\varepsilon$. We determine $t_c$ from the shape
of this $n_L(\varepsilon)$ curve. A charge sensor with different
responses to occupation of the two dots lets us infer $n_L$
from its mean signal.

\subsection{Equilibrium charge distribution}
\label{sec:measurement-equilibrium-charge}
Keep one orbital--valley state per dot and a fixed spin state.
The localized charge states are $\ket L=\ket{1,0}$ and
$\ket R=\ket{0,1}$. Removing their common mean energy and
choosing the relative basis phase gives
\begin{equation}
 H_c=\begin{pmatrix}\varepsilon/2&-t_c\\-t_c&-\varepsilon/2\end{pmatrix},
 \qquad
 E_\pm=\pm\frac{\Delta_c}{2},\quad
 \Delta_c=\sqrt{\varepsilon^2+4t_c^2}.
 \label{eq:measurement-charge-anticrossing}
\end{equation}
For $t_c>0$, define $0<\theta<\pi$ by
$\cos\theta=\varepsilon/\Delta_c$ and
$\sin\theta=2t_c/\Delta_c$. The normalized eigenstates are
\begin{equation}
 \ket g=\sin\frac\theta2\ket L+\cos\frac\theta2\ket R,
 \qquad
 \ket e=\cos\frac\theta2\ket L-\sin\frac\theta2\ket R,
 \label{eq:measurement-charge-eigenstates}
\end{equation}
with energies $E_-$ and $E_+$, respectively. At zero temperature
the electron occupies $\ket g$, so the mean left-dot occupation is
\begin{equation}
 n_L^{(g)}=|\langle L|g\rangle|^2
       =\frac12\left(1-\frac{\varepsilon}
                            {\sqrt{\varepsilon^2+4t_c^2}}\right).
 \label{eq:measurement-ground-charge-occupation}
\end{equation}
It approaches one for $\varepsilon\ll-2t_c$, one half at zero,
and zero for $\varepsilon\gg2t_c$. Its central slope is
$dn_L^{(g)}/d\varepsilon|_0=-1/(4t_c)$: increasing $t_c$ makes
the fall in $n_L^{(g)}$ less steep as a function of $\varepsilon$,
quantifying the broader transition described above.

Finite temperature also broadens this curve by populating
$\ket e$, whose left-dot occupation is
$n_L^{(e)}=(1+\varepsilon/\Delta_c)/2$.
If the two levels reach thermal equilibrium at temperature $T_e$,
then $p_e/p_g=e^{-\Delta_c/(k_BT_e)}$ and $p_e+p_g=1$.
Taking $n_L=p_g n_L^{(g)}+p_e n_L^{(e)}$ gives
\begin{equation}
 n_L\equiv\langle\widehat N_1\rangle
    =\frac12\left[1-\frac{\varepsilon}{\Delta_c}
                  \tanh\!\left(\frac{\Delta_c}{2k_BT_e}\right)\right].
 \label{eq:measurement-charge-polarization}
\end{equation}
For $t_c\ll k_BT_e$, this approaches
$n_L=1/[1+e^{\varepsilon/(k_BT_e)}]$, which hardly depends on
$t_c$. A thermally dominated curve cannot accurately resolve a
much smaller coupling even when the temperature is known.

\subsection{From the sensor record to the coupling}
Because the total charge remains one,
$\langle\widehat N_2\rangle=1-n_L$. Substitution into the
sensor model of Eq.~\eqref{eq:measurement-sensor-signal} gives
\begin{equation}
 \overline s=s_{\mathrm{bg}}+\kappa_2
                         +(\kappa_1-\kappa_2)n_L.
 \label{eq:measurement-interdot-sensor-signal}
\end{equation}
The bar denotes an average of repeated or stationary records.
The sensor must have $\kappa_1\ne\kappa_2$ to detect redistribution
at fixed total charge. After subtracting the smooth gate-dependent
background, call the signal $s'$. Its values far on the left- and
right-occupied sides are $s_L=\kappa_1$ and $s_R=\kappa_2$.
The normalized measured curve is then
\begin{equation}
 n_L=\frac{s'-s_R}{s_L-s_R}.
 \label{eq:measurement-normalized-charge-signal}
\end{equation}

Load one electron and close its barriers to the reservoirs.
Keep the interdot barrier setting fixed and scan a chosen plunger
voltage $v$ across the redistribution region. Allow the two-level
system to relax at each setting before averaging the sensor signal.
Use a scan range in which $t_c$ is approximately constant and
the detuning is approximately linear,
$\varepsilon=\eta(v-v_0)$; here $v_0$ is the equal-energy setting
and $\eta=\partial\varepsilon/\partial v$ has units of energy
per voltage. The common shift of the two dot energies does not
affect this fixed-charge model.

Record both nearly localized ends of the scan. There the dot
occupation is nearly constant, so the remaining signal slope
comes from the direct gate response. For a narrow scan with an
approximately linear background, fit the two ends with a common
slope and separate offsets, then subtract that common slope.
The two corrected plateaus supply $s_L$ and $s_R$; any common
constant offset cancels in the normalization. Normalize using
Eq.~\eqref{eq:measurement-normalized-charge-signal}, then compare
the measured curve with Eq.~\eqref{eq:measurement-charge-polarization}.
For independently known $\eta$ and $T_e$, the remaining shape
parameter is $t_c$; $v_0$ sets the center. In the low-temperature
limit the central slope gives directly
\[
 t_c=\frac{|\eta|}{4|dn_L/dv|_{v_0}},
\]
while fitting the full curve uses more of the data. This is a
measurement of coupling through the equilibrium charge distribution,
without transport through the double dot~\cite{DiCarlo2004}.

\subsection{Spectroscopy with an oscillating gate}
\label{sec:measurement-charge-drive-spectroscopy}
There is also a frequency-based measurement. An oscillating gate
voltage changes the relative energies of the two dots. Because
$\ket g$ and $\ket e$ are superpositions of the localized states,
this perturbation can drive transitions between them. Populating
$\ket e$ changes the mean charge distribution, which the same
sensor can detect~\cite{Petta2004Charge}.

Suppose the drive changes detuning by
$A_\varepsilon\cos(2\pi ft)$ while leaving $t_c$ approximately
fixed. With $\tau_z=\ket L\bra L-\ket R\bra R$,
\begin{equation}
 H_d(t)=\frac{A_\varepsilon}{2}\cos(2\pi ft)\tau_z,
 \qquad
 \langle e|\tau_z|g\rangle=\sin\theta=\frac{2t_c}{\Delta_c}.
 \label{eq:measurement-charge-drive}
\end{equation}
Thus the transition matrix element is nonzero. The weak-drive
resonance condition is the same two-level condition derived in
Section~\ref{sec:transverse-spin-control}, now with gap $\Delta_c$:
\begin{equation}
 hf=\Delta_c(\varepsilon),\qquad
 \varepsilon_{\mathrm{res}}
      =\pm\sqrt{(hf)^2-4t_c^2},\quad hf\geq2t_c.
 \label{eq:measurement-photon-assisted-tunneling}
\end{equation}
Take $k_BT_e\ll\Delta_c$, so the drive-off state is approximately
$\ket g$. Under a weak continuous drive, excitation competes with
relaxation and produces a mean excited-state population $p_e$.
A sensor averaging over many drive periods measures the population
contribution; the rapidly oscillating coherence averages away.
Since $n_L^{(e)}-n_L^{(g)}=\varepsilon/\Delta_c$, the change is
\begin{equation}
 \delta\overline s=(\kappa_1-\kappa_2)p_e
                         \frac{\varepsilon}{\Delta_c}.
 \label{eq:measurement-charge-resonance-signal}
\end{equation}
The resonantly increased population therefore gives a sensor
feature away from zero detuning. At zero detuning both eigenstates
have equal mean dot charge, so this particular signal vanishes
even though the drive matrix element is nonzero.

At a fixed weak drive amplitude and frequency $f$, scan $v$ and
subtract the drive-off sensor record from the drive-on record.
Locate the narrow features on the two sides of $v_0$, and repeat
for several frequencies. Their resonance positions obey
\begin{equation}
 (hf)^2=\eta^2(v_{\mathrm{res}}-v_0)^2+4t_c^2.
 \label{eq:measurement-charge-resonance-fit}
\end{equation}
For a symmetric local response, the midpoint of each pair gives
$v_0$. Plotting $(hf)^2$ against $(v_{\mathrm{res}}-v_0)^2$
gives slope $\eta^2$ and intercept $4t_c^2$. This determines both
$|\eta|$ and $t_c$, even when the feature at the minimum gap is
invisible. Use resolved, weak-drive resonances for which the
charge-contrast factor varies little across each linewidth;
otherwise a signal extremum need not coincide with the resonance.
